\chapter{Review of category theory}
\label{chap:constrcat}

This chapter recalls classic constructions of category theory. We discuss over and under categories, inverse image categories, and diagram categories. The construction of limits and colimits is recalled. We discuss cofinal functors, the Grothendieck construction and some of its basic properties.

\section{Basic definitions and notations}
\label{sec:basicdefnotations}

\subsection{Initial and terminal categories}
\label{sec:initterminalcat}
The initial object in a category is denoted $\initial$, and the terminal object $\terminal$. The initial category (with an empty set of objects) is denoted $\initialcat$, and the terminal category (with one object and one identity map) is denoted $\terminalcat$.\index{$\initial$, $\terminal$, $\initialcat$, $\terminalcat$}

For a category ${\cat D}$ we denote $e_d \colon \terminalcat \ra {\cat D}$\index{$e_d$, $p_{\cat D}$} the functor that embeds $\terminalcat$ as the object $d \eps {\cat D}$, and $p_{\cat D} \colon {\cat D} \ra \terminalcat$ the terminal category projection.

\subsection{Inverse image category}
\label{sec:catinverseimage}
\index{category!inverse image category}
For a functor $u \colon {\cat A} \ra {\cat B}$ and an object $b$ of ${\cat B}$, the {\it inverse image} of $b$ is the subcategory $u^{-1}b$ of ${\cat A}$ consisting of objects $a$ with $ua = b$ and maps $f \colon a \ra a^{'}$ with $u f = 1_b$.

\subsection{Categories of diagrams}
\label{sec:catdiag}
If ${\cat D}$ and ${\cat M}$ are categories, the category of functors ${\cat D} \ra {\cat M}$ (or ${\cat D}$-diagrams of ${\cat M}$) is denoted ${\cat M}^{\cat D}$. 

A functor $u \colon {\cat D}_1 \ra {\cat D}_2$ induces a functor of diagram categories denoted $u^* \colon {\cat M}^{{\cat D}_2} \ra {\cat M}^{{\cat D}_1}$. If two functors $u, v$ are composable, then $(uv)^* = v^*u^*$. If two functors $u \dashv v$ are adjoint, then $v^* \dashv u^*$ are adjoint.

If ${\cat C}$ is a class of maps of ${\cat M}$ (for example the weak equivalences or the cofibrations in a cofibration category), we denote ${\cat C}^{\cat D}$ the class of maps $f$ of ${\cat M}^{\cat D}$ such that $f_d \eps C$ for any object $d \eps {\cat D}$.

We'd like to point out that we have defined the weak equivalences ${\cat W}$ and the cofibrations ${\cat Cof}$ of a cofibration category $({\cat M}$, ${\cat W}$, ${\cat Cof})$ as {\it classes of maps} rather than as subcategories. For a cofibration category, ${\cat W}^{\cat D}$ (and ${\cat Cof}^{\cat D}$) therefore denote the class of diagram maps $f \colon X \ra Y$ in ${\cat M}^{\cat D}$ such that each $f_d$ is a weak equivalence, resp. a cofibration. If $w{\cat M}$ denotes the subcategory of ${\cat M}$ generated by ${\cat W}$, then $(w{\cat M})^{\cat D}$ is {\it not} the same thing as ${\cat W}^{\cat D}$.

\section{Limits and colimits}
\label{sec:relativelimcolim}
Assume ${\cat D}_1$ and ${\cat D}_2$ are categories, $u \colon {\cat D}_1 \ra {\cat D}_2$ is a functor and ${\cat M}$ is a category. A colimit functor along $u$, denoted $\colim^u \colon {\cat M}^{{\cat D}_1} \ra {\cat M}^{{\cat D}_2}$, is by definition a left adjoint of $u^* \colon {\cat M}^{{\cat D}_2} \ra {\cat M}^{{\cat D}_1}$. A limit functor along $u$, denoted $\lim^u \colon {\cat M}^{{\cat D}_1} \ra {\cat M}^{{\cat D}_2}$, is by definition a right adjoint of $u^*$. $\colim^u$ is also called the left Kan extension along $u$, and $\lim^u$ is called the right Kan extension along $u$. 

Being a left adjoint, if $\colim^u$ exists then it is unique up to a unique natural isomorphism. As a right adjoint, if $\lim^u$ exists it is unique up to a unique natural isomorphism. 

We will also consider the case when $\colim^uX$ (resp. $\lim^uX$) exists for some, but not all objects $X$ of ${\cat M}^{{\cat D}_1}$. Each of $\colim^uX$ and $\lim^uX$ are defined by an universal property, and if $\colim^uX$ or $\lim^uX$ exist, then they are unique up to unique isomorphism.

If ${\cat D}_1 = {\cat D}$ and ${\cat D}_2$ is the point category, $\colim^u$ is the well-known 'absolute' colimit of ${\cat D}$-diagrams, denoted $\colim^{\cat D}$ (or simply $\colim$ if there is no confusion). Dually, in this case $\lim^u$ is the 'absolute' limit $\lim^{\cat D}$. If ${\cat D}$ has a terminal object $\final$, then $\colim^{\cat D} X$ always exists and the natural map $X_{\final} \ra \colim^{\cat D} X$ is an isomorphism. If ${\cat D}$ has an initial object $\initial$ then $\lim^{\cat D} X$ always exists and the natural map $\lim^{\cat D} X \ra X_{\initial}$ is an isomorphism.

Next lemma presents a base change formula for relative (co)limits. In the case of colimits, the lemma states that the relative colimit along a functor $u \colon {\cat D}_1 \ra {\cat D}_2$ can be pointwisely computed in terms of absolute colimits over $(u \downarrow d_2)$ for all objects $d_2 \eps {\cat D}_2$.

\begin{lem}
\label{lem::computerelcolim}
Suppose that $u \colon {\cat D}_1 \ra {\cat D}_2$ is a functor, and suppose that $X$ is a ${\cat D}_1$ diagram in a category ${\cat M}$.
\begin{enumerate}
\item
If $\colim^{(u \downarrow d_2)} X$ exists for all $d_2 \eps {\cat D}_2$, then $\colim^u X$ exists and 
\begin{center}
${\colim}^{(u \downarrow d_2)} X \cong ({\colim}^u X)_{d_2}$
\end{center}
\item
If $\lim^{(d_2 \downarrow u)} X$ exists for all $d_2 \eps {\cat D}_2$, then $\lim^u X$ exists and 
\begin{center}
$({\lim}^u X)_{d_2} \cong {\lim}^{(d_2 \downarrow u)} X$
\end{center}
\end{enumerate}
\end{lem}

\begin{proof}
We only prove (1). We show that $({\colim}^{(u \downarrow d_2)} X)_{d_2 \eps {\cat D}_2}$ satisfies the universal property that defines $\colim^u X$. Maps from $X$ to a diagram $u^* Y$ can be identified with maps $X_{d_1} \ra Y_{ud_1}$ that make the diagram below commutative for all maps $f \colon d_1 \ra d_1^{'}$
\begin{center}
$\xymatrix{
    X_{d_1} \ar[r] \ar[d]_{X_f} &  Y_{ud_1} \ar[d]^{Y_{uf}} \\
    X_{d_1^{'}} \ar[r] & Y_{ud_1^{'}}
    }$
\end{center}
They can be further identified with maps $X_{d_1} \ra Y_{d_2}$, defined for all $\phi \colon ud_1 \ra d_2$, that make the diagram below commutative 
\begin{center}
$\xymatrix{
    X_{d_1} \ar[r] \ar[d]_{X_f} &  Y_{d_2} \ar[d]^{Y_g} \\
    X_{d_1^{'}} \ar[r] & Y_{d_2^{'}}
    }$
\end{center}
for all maps $f, g$ that satisfy $\phi^{'} uf = g \phi$. They can  finally be identified with maps ${\colim}^{(u \downarrow d_2)} X \ra Y_{d_2}$ compatible in $d_2$, and from the universal property of the colimit we see that ${\colim}^u X$ exists and ${\colim}^{(u \downarrow d_2)} X \cong ({\colim}^u X)_{d_2}$.
\end{proof}

In particular, if ${\cat M}$ is closed under absolute colimits, it is also closed under relative colimits. By definition, if ${\cat M}$ is closed under small colimits we say that ${\cat M}$ is {\it cocomplete}. Dually, if ${\cat M}$ is closed under absolute limits, it is also closed under relative limits. If ${\cat M}$ is closed under small limits we say by definition that ${\cat M}$ is {\it complete}.

The following lemma presents another well known result - that the composition of (co)limits is the (co)limit of the composition.


\begin{lem}
\label{lem::colimcomposition}
Suppose that $u \colon {\cat D}_1 \ra {\cat D}_2$ and $v \colon {\cat D}_2 \ra {\cat D}_3$ are two functors, and suppose that $X$ is a ${\cat D}_1$ diagram in a category ${\cat M}$.
\begin{enumerate}
\item
Assume that $\colim^u X$ exists. If either $\colim^v \colim^u X$ or \\ $\colim^{vu} X$ exist, then they both exist and are canonically isomorphic.
\item
Assume that $\lim^u X$ exists. If either $\lim^v \lim^u X$ or $\lim^{vu} X$ exist, then they both exist and are canonically isomorphic.
\end{enumerate}
\end{lem}

\begin{proof}
We only prove (1). Since $\colim^u X$ exists, we have a bijection of sets natural in $Y \eps {\cat M}^{{\cat D}_3}$
\begin{center}
$Hom(X, u^*v^* Y) \cong Hom(\colim^u X, v^* Y)$ 
\end{center}
If $\colim^{vu} X$ exists, we also have a natural bijection
\begin{center}
$Hom(X, u^*v^* Y) \cong Hom(\colim^{vu} X, Y)$
\end{center}
therefore $\colim^{vu} X$ satisfies the universal property of $\colim^v \colim^u X$. If on the other hand $\colim^v \colim^u X$ exists, we have a natural bijection
\begin{center}
$Hom(\colim^u X, v^* Y) \cong Hom(\colim^v\colim^u X, Y)$
\end{center}
and $\colim^v \colim^u X$ satisfies the universal property of $\colim^{vu} X$. 
\end{proof}

In particular, if ${\cat M}$ is cocomplete then $\colim^{vu}$ and $\colim^v \colim^u$ are naturally isomorphic. Dually, if ${\cat M}$ is complete then $\lim^{vu}$ and $\lim^v \lim^u$ are naturally isomorphic.

\section{Simplicial sets}
\label{sec:simplindex}
\index{simplicial set}
This text assumes familiarity with simplicial sets, and the reader may refer for example to \cite{Goerss-Jardine} or \cite{Hovey} for a treatment of the standard theory of simplicial sets. 

We will denote by ${\bf n}$ the poset $\{0, 1, ... , n\}$, for $n \ge 0$, with the natural order. The cosimplicial indexing category $\Delta$\index{$\Delta$} is the category with objects ${\bf 0}, {\bf 1}, ..., {\bf n}, ...$ and maps the order-preserving maps ${\bf n}_1 \ra {\bf n}_2$. If ${\cat C}$ is a category, a simplicial object in ${\cat C}$ is then a functor $\Delta^{op} \ra {\cat C}$, and a cosimplicial object in ${\cat C}$ is a functor $\Delta \ra {\cat C}$.

\comment{
The maps of $\Delta$ are generated by the maps below
\begin{enumerate}
\item[]
$d^i \colon {\bf n-1} \ra {\bf n}$ for $0 \le i \le n, \; n \ge 1$ (the coface maps)
\begin{align*}
d^i(k) = k \;\;\;\;\;\;\;\;\;\; (k < i) \\
d^i(k) = k+1 \;\;\;\; (k \ge i)
\end{align*}
\item[]
$s^i \colon {\bf n+1} \ra {\bf n}$ for $0 \le i \le n, \; n \ge 0$ (the codegeneracy maps )
\begin{align*}
s^i(k) = k \;\;\;\;\;\;\;\;\;\; (k \le i) \\
s^i(k) = k-1 \;\;\;\; (k > i)
\end{align*}
\end{enumerate}
subject to the relations given by the {\it cosimplicial identities}
\begin{align}
d^jd^i = d^id^{j-1} \;\;\;\;\;\;\, (i < j) \;\;\;\;\;\, \notag \\
s^jd^i = d^is^{j-1} \;\;\;\;\;\;\, (i < j) \;\;\;\;\;\, \notag \\
s^jd^j = s^jd^{j+1} = 1  \;\;\;\;\;\;\;\;\;\;\;\;\;\;\;\,  \label{eqn:cosimplidentities} \\
s^jd^i = d^{i-1}s^j \;\;\;\;\;\;\, (i > j+1) \notag \\
s^js^i = s^is^{j+1} \;\;\;\;\;\;\, (i \le j) \;\;\;\;\;\, \notag
\end{align}

Therefore, in order to define a cosimplicial object $X\tdot$ in ${\cat C}$ it suffices to define the objects $X^n, n \ge 0$ and maps
\begin{align*}
d^i \colon X^{n-1} \ra  X^n \;\;\;\;\;\;\, (0 \le i \le n, \; n \ge 1) \\
s^i \colon X^{n+1} \ra X^{n} \;\;\;\;\;\;\, (0 \le i \le n, \; n \ge 0)
\end{align*}
satisfying the cosimplicial identities (\ref{eqn:cosimplidentities}). Dually, in order to define a simplicial object $X\sdot$ in ${\cat C}$ it suffices to define the objects $X_n, n \ge 0$, the face maps $d_i$ and the degeneracy maps $s_i$
\begin{align*}
d_i \colon X_n \ra X_{n-1} \;\;\;\;\;\;\, (0 \le i \le n, \; n \ge 1) \\
s_i \colon X_n \ra X_{n+1} \;\;\;\;\;\;\, (0 \le i \le n, \; n \ge 0)
\end{align*}
subject to the {\it simplicial identities} relations dual to the cosimplicial identities.
\begin{align}
d_id_j = d_{j-1}d_i \;\;\;\;\;\;\, (i < j) \notag \;\;\;\;\;\, \\
d_is_j = s_{j-1}d_i \;\;\;\;\;\;\, (i < j) \notag \;\;\;\;\;\, \\
d_js_j = d_{j+1}s_j = 1  \;\;\;\;\;\;\;\;\;\;\;\;\;\;\;\,  \label{eqn:simplidentities}\\
d_is_j = s_jd_{i-1} \;\;\;\;\;\;\,  (i > j+1) \notag \\
s_is_j = s_{j+1}s_i \;\;\;\;\;\;\, (i \le j) \;\;\;\;\;\, \notag
\end{align}

A cosimplicial object map $f \colon X\tdot \ra Y\tdot$ is a natural transformation of the corresponding functors $\Delta \ra {\cat C}$, or equivalently is a collection of maps $X^n \ra Y^n$ that commute with the coface and codegeneracy maps. Dually, a simplicial object map $f \colon X\sdot \ra Y\sdot$ can be viewed as a collection of maps $X_n \ra Y_n$ that commute with the face and degeneracy maps.
} % end comment

Taking in particular simplicial objects in $Sets$ we get the category of simplicial sets, denoted $sSets$. The representable functor ${\Delta^{op}}(-, {\bf n})$ determines a simplicial set denoted $\Delta[n]$, called the standard $n$-simplex (for $n \ge 0$). For any map ${\bf m} \ra {\bf n}$ we have a simplicial set map $\Delta[m] \ra \Delta[n]$, thus $\Delta[-]$ determines a cosimplicial object in $sSets$.

In particular, we have simplicial set inclusions $d^0, d^1 \colon \Delta[0] \ra \Delta[1]$, which induce respectively functors $i_0, i_1 \colon X\sdot \cong X\sdot \times \Delta[0] \ra  X\sdot \times \Delta[1]$. We say that two maps $f, g \colon X\sdot \ra Y\sdot$ are simplcially homotopic if there exists a simplicial set map $h \colon X\sdot \times \Delta[1] \ra Y\sdot$ such that $f = hi_0, g = hi_1$. 

%Since $\Delta[1]_n = {\Delta}({\bf n}, {\bf 1})$, a simplicial homotopy $h \colon X\sdot \times \Delta[1] \ra Y\sdot$ is specified by maps $h_i \colon X_n \ra Y_n$, for $0 \le i \le n+1, n \ge 0$, such that $h_0 = g$ and $h_{n+1} = f$ in degree $n$ (notice the reversal!) and 
%\begin{align}
%d_ih_j = h_{j-1}d_i \;\;\;\, (i < j) \notag \\
%d_ih_j = h_jd_i     \;\;\;\;\;\;\, (i \ge j)  \label{eqn:simplhomotidentities} \\
%s_ih_j = h_{j+1}s_i \;\;\;\, (i < j) \notag \\
%s_ih_j = h_js_i     \;\;\;\;\;\;\, (i \ge j) \notag
%\end{align}

%The identities (\ref{eqn:simplhomotidentities}) allow us to define simplicial homotopy for maps between simplicial objects in an arbitrary category. Dualizing the identities (\ref{eqn:simplhomotidentities}), we can define cosimplicial homotopy for maps between cosimplicial objects.

%Equivalently, a simplicial homotopy between $f, g \colon X\sdot \ra Y\sdot$ can be specified by maps $h^{'}_i = h_id_i \colon X_{n+1} \ra Y_n$, for $0 \le i \le n+1, n \ge 0$, such that $h^{'}_0s_0 = g$ and $h^{'}_{n+1}s_n = f$ in degree $n$ and
%\begin{align}
%d_ih^{'}_j = h^{'}_{j-1}d_i \;\;\;\, (i < j) \notag \\
%d_ih^{'}_j = h^{'}_jd_{i+1} \;\;\;\, (i \ge j)  \label{eqn:simplhomotidentities2} \\
%s_ih^{'}_j = h^{'}_{j+1}s_i \;\;\;\, (i < j) \notag \\
%s_ih^{'}_j = h^{'}_js_{i+1} \;\;\;\, (i \ge j) \notag
%\end{align}

%The maps $h_i$ are recovered from $h^{'}_i$ using in degree $n$ the relation $h_i = h^{'}_is_i$ for $i \le n$  and $h_i = h^{'}_is_{i-1}$ for $i > 0$.

In general, the simplicial homotopy relation is not reflexive nor transitive.

%but we may consider the equivalence relation it generates. 
%We say that a map $f \colon X\sdot \ra Y\sdot$ is a {\it simplicial homotopy equivalence} if there is a map in the opposite direction $g \colon Y\sdot \ra X\sdot$ such that $gf, 1_{X\sdot}$ resp. $fg, 1_{Y\sdot}$ are in the same simplicial homotopy relation equivalence class.

\section{The nerve of a category}
\label{sec:catnerve}
Given a category ${\cat C}$, its nerve (or classifying space) $B{\cat C}$ is a simplicial set having as $n$-simplices the composable strings $A_0 \ra A_1 \ra ... \ra A_n$ of maps in ${\cat C}$. The poset ${\bf n}$ can be viewed as a category with objects $\{0, 1, ..., n\}$ and maps $n_1 \ra n_2$ for all $n_1 \le n_2$. $B{\cat C}_n$ can then be identified with the set of functors ${\bf n} \ra {\cat C}$, and we define the maps $({\bf n_1} \ra {\cat C}) \lra ({\bf n_2} \ra {\cat C})$ to be the commutative diagrams
\begin{center}
$ \xymatrix {
        {\bf n_2} \ar[rr] \ar[dr] & & {\bf n_1} \ar[dl] \\
	& {\cat C} &
	}$
\end{center}
The nerve functor $B{\cat C}$ commutes with arbitrary limits (has a left adjoint).

Denote $I$ the category ${\bf 1}$, with two objects and only one non-identity map. Its nerve $BI$ is isomorphic to $\Delta[1]$.

A functor $f \colon {\cat C}_1 \ra {\cat C}_2$ induces a map of simplicial sets $Bf \colon B{\cat C}_1 \ra B{\cat C}_2$. A natural map between two functors $h \colon f \Ra g$ induces a functor $H \colon {\cat C}_1 \times I \ra {\cat C}_2$, therefore a simplicial homotopy $BH \colon B{\cat C}_1 \times \Delta[1] \ra B{\cat C}_2$ between $Bf$ and $Bg$.

\bigskip

\section{Cofinal functors}
\label{sec:coffunctors1}
This section is a short primer on cofinal and homotopy cofinal functors. We only prove the minimal set of properties that we will need for the rest of the text. For more information on cofinal and homotopy cofinal functors, please refer to Hirschhorn \cite{Hirschhorn}.

\index{functor!left, right cofinal}
\index{functor!homotopy left, right cofinal}
\begin{defn}
\label{defn:coffunct}
A functor $u \colon {\cat D}_1 \ra {\cat D}_2$ beween small categories is
\begin{enumerate}
\item
{\it left cofinal} if for any object $d_2$ of ${\cat D}_2$ the space $B(u \downarrow d_2)$ is connected and non-empty,
\item
{\it homotopy left cofinal} if for any object $d_2$ of ${\cat D}_2$ the space $B(u \downarrow d_2)$ is contractible,
\item
{\it right cofinal} if for any object $d_2$ of ${\cat D}_2$ the space $B(d_2 \downarrow u)$ is connected and non-empty
\item
{\it homotopy right cofinal} if for any object $d_2$ of ${\cat D}_2$ the space $B(d_2 \downarrow u)$ is contractible
\end{enumerate}
\end{defn}

A homotopy left (resp. right) cofinal functor is in particular left (resp. right) cofinal. 

\begin{lem}
\label{lem::adjfunctcofinalfunct}
If $u : {\cat D}_1 \rightleftarrows {\cat D}_2 : v$ is an adjoint functor pair, then $u$ is homotopy left cofinal and $v$ is homotopy right cofinal.
\end{lem}
\begin{proof}
For any object $d_2$ of ${\cat D}_2$, the under category $(u \downarrow d_2)$ is isomorphic to $({\cat D}_1 \downarrow vd_2)$, and the latter has $id \colon vd_2 \ra vd_2$ as a terminal object. $B(u \downarrow d_2)$ is therefore contractible for any $d_2$, so $u$ is homotopy left cofinal.

A dual proof shows that $v$ is homotopy right cofinal.
\end{proof}

The following lemma gives a characterization of left (resp. right) cofinal functors in terms of their preservation of limits (resp. colimits).

\begin{prop}
\label{prop:cofinallimitconv}
Suppose that $u \colon {\cat D}_1 \ra {\cat D}_2$ is a functor of small categories. Then:
\begin{enumerate}
\item
$u$ is right cofinal iff for any cocomplete category ${\cat M}$ and any diagram $X \eps {\cat M}^{{\cat D}_2}$ the natural map $\colim^{{\cat D}_1}u^*X \ra \colim^{{\cat D}_2}X$ is an isomorphism in ${\cat M}$.
\item
$u$ is left cofinal iff for any complete category ${\cat M}$ and any diagram $X \eps {\cat M}^{{\cat D}_2}$ the natural map $\lim^{{\cat D}_2}X \ra \lim^{{\cat D}_1}u^*X$ is an isomorphism in ${\cat M}$.
\end{enumerate}
\end{prop}

Note the subtle positional difference of $u$ between the formulas of \lemmaref{lem::computerelcolim} (1) and \propositionref{prop:cofinallimitconv} (2). For a diagram $Y \eps {\cat M}^{{\cat D}_1}$ and an object $d_2 \eps {\cat D}_2$ we have that $(\colim^u Y)_{d_2} \cong \colim^{(u \downarrow d_2)}Y$. But for a diagram $X \eps {\cat M}^{{\cat D}_2}$, the natural map $\colim^{{\cat D}_1}u^*X \ra \colim^{{\cat D}_2}X$ is an isomorphism iff $(d_2 \downarrow u)$ is connected and non-empty for all objects $d_2 \eps {\cat D}_2$.

\begin{proof}[Proof of \propositionref{prop:cofinallimitconv}]
\mbox{}

(1) ($\La$). Let ${\cat M}$ be ${\cat Set}$ and $X$ be the ${\cat D}_2$-diagram ${{\cat D}_2}(d_2, -)$ for $d_2 \eps {\cat D}_2$. We notice that $\colim^{{\cat D}_2}X \cong \colim^{d^{'}_2 \eps {\cat D}_2} {{\cat D}_2}(d_2, d^{'}_2)$ is the one-point set, and $\colim^{{\cat D}_1}u^*X \cong \colim^{d_1 \eps {\cat D}_1} {{\cat D}_2}(d_2, ud_1)$ is the set of connected components of $B(d_2 \downarrow u)$, which must therefore be isomorphic with the one-point set. The conclusion is now proved.

(1) ($\Ra$). This will be a consequence of \propositionref{prop:cofinallimit} below. 

The proof of (2) is dual to the proof of (1).
\end{proof}

We next prove something a bit stronger than the right to left implication of \propositionref{prop:cofinallimitconv} (1).

\begin{prop}
\label{prop:cofinallimit}
Suppose that $u \colon {\cat D}_1 \ra {\cat D}_2$ is a functor, that ${\cat M}$ is a category and that $X$ an object of ${\cat M}^{{\cat D}_2}$.
\begin{enumerate}
\item
Assume that the functor $u$ is right cofinal. If either $\colim^{{\cat D}_1}u^*X$ or $\colim^{{\cat D}_2}X$ exist, then they both exist and the natural map $\colim^{{\cat D}_1}u^*X \ra \colim^{{\cat D}_2}X$ is an isomorphism.
\item
Assume that the functor $u$ is left cofinal. If either $\lim^{{\cat D}_1}u^*X$ or $\lim^{{\cat D}_2}X$ exist, then they both exist and the natural map $\lim^{{\cat D}_2}X \ra \lim^{{\cat D}_1}u^*X$ is an isomorphism.
\end{enumerate}
\end{prop}

\begin{proof}
We will prove (1); the statement (2) follows from duality. We denote $p_{{\cat D}_i} \colon {\cat D}_i \ra \terminalcat$, $i = 1, 2$ the terminal category projections.

We start by constructing a natural isomorphism in $X \eps {\cat M}^{{\cat D}_2}, Y \eps {\cat M}$
\begin{equation}
\label{eqn:cofinalcolimitiso}
{{\cat M}^{{\cat D}_1}}(u^*X, p_{{\cat D}_1}^*Y) \cong {{\cat M}^{{\cat D}_2}}(X, p_{{\cat D}_2}^*Y)
\end{equation}
From right to left, if $f \colon X \ra p_{{\cat D}_2}^*Y \eps {\cat M}^{{\cat D}_2}$ then we define $F(f) \colon u^*X \ra p_{{\cat D}_1}^*Y \eps {\cat M}^{{\cat D}_1}$ on component $d_1 \eps {\cat D}_1$ as $F(f)_{d_1} = f_{ud_1} \colon X_{ud_1} \ra Y$.

From left to right, if $f \colon u^*X \ra p_{{\cat D}_1}^*Y \eps {\cat M}^{{\cat D}_1}$ then we define $G(f) \colon X \ra p_{{\cat D}_2}^*Y \eps {\cat M}^{{\cat D}_2}$ on component $d_2 \eps {\cat D}_2$ as follows. Since $(d_2 \downarrow u)$ is non-empty we can pick a map $\alpha \colon d_2 \ra ud_1$, and we define $G(f)_{d_2} \colon X_{d_2} \ra X_{ud_1} \ra Y$ to be the composition of $X_\alpha$ with $f_{d1}$. Since $(d_2 \downarrow u)$ is connected, the map $G(f)$ does not depend on the choice involved, $G(f)$ is indeed a diagram map from $X$ to $p_{{\cat D}_2}^*Y$ and furthermore $F, G$ are inverses of each other.

If we view both terms of \ref{eqn:cofinalcolimitiso} as functors of $Y$, then $\colim^{{\cat D}_1}u^*X$ exists iff the left side of \ref{eqn:cofinalcolimitiso} is a representable functor of $Y$ iff the right side of \ref{eqn:cofinalcolimitiso} is a representable functor of $Y$ iff $\colim^{{\cat D}_2}X$ exists. Under these conditions it also follows that the natural map $\colim^{{\cat D}_1}u^*X \ra \colim^{{\cat D}_2}X$ is an isomorphism.
\end{proof}

As a consequence of \propositionref{prop:cofinallimitconv}, it is not hard to see that left (resp. right) cofinal functors are stable under composition. 

\section{The Grothendieck construction}
\label{sec:grothendieckconstruction}
\index{Grothendieck construction}
Denote ${\cat Cat}$ the category of small categories and functors. If ${\cat D}$ is a small category and $H$ is a functor ${\cat D}$ $\ra$ ${\cat Cat}$, the Grothendieck construction of $H$, denoted $\int_{\cat D} H$, is the category for which:
\begin{enumerate}
\item
objects are pairs $(d, x)$ with $d$ an object of ${\cat D}$ and $x$ an object of the category $H(d)$
\item 
maps $(d_1, x_1) \ra (d_2, x_2)$ are pairs of maps $(f, \phi)$ with $f \colon d_1 \ra d_2$ and $\phi \colon H(f)x_1 \ra x_2$
\item 
the composition of maps $(g, \psi)(f, \phi)$ is the map $(gf,  \psi \comp H(g) \phi)$
\item 
the identity of $(d, x)$ is $(1_d, 1_x)$
\end{enumerate}

The Grothendieck construction comes with a projection functor $p \colon \int_{\cat D} H \ra {\cat D}$, defined by $(d, x) \ra d$. If $d$ is an object of ${\cat D}$, then the inverse image category $p^{-1}d$ may be identified with the category $H(d)$. 

Dually if $H$ is a functor ${\cat D}^{op} \ra {\cat Cat}$, the contravariant Grothendieck construction of $H$ is defined as $\int^{\cat D} H = (\int_{{\cat D}^{op}} {H^{op}})^{op}$, and comes with a projection functor $p \colon \int^{\cat D} H \ra {\cat D}$.

\begin{prop}
\label{prop:grothendieckconstructionadjoint}
\mbox{}
\begin{enumerate}
\item
Suppose that $H \colon {\cat D} \ra {\cat Cat}$ is a functor from a small category to the category of small categories, with projection functor $p \colon \int_{\cat D} H \ra {\cat D}$. Then for any object $d \eps {\cat D}$, the inclusion $p^{-1}d \ra (p \downarrow d)$ admits a left adjoint.
\item
Suppose that $H \colon {\cat D}^{op} \ra {\cat Cat}$ is a functor from a small category to the category of small categories, with projection functor $p \colon \int^{\cat D} H \ra {\cat D}$. Then for any object $d \eps {\cat D}$, the inclusion $p^{-1}d \ra (d \downarrow p)$ admits a right adjoint.
\end{enumerate}
\end{prop}

\begin{proof}
To prove (1), denote $i_d \colon p^{-1}d \ra (p \downarrow d)$ the inclusion. A left adjoint $j_d \colon (p \downarrow d) \ra p^{-1}d$ of $i_d$ can be constructed such that:
\begin{enumerate}
\item[(a)]
$j_d$ sends the object $((d_1, x_1), f_1 \colon d_1 \ra d)$ of $(p \downarrow d)$ to the object $H(f_1)x_1$ of $H(d) \cong p^{-1}d$
\item[(b)]
For objects $((d_1, x_1), f_1 \colon d_1 \ra d)$ and $((d_2, x_2), f_2 \colon d_2 \ra d)$ of $(p \downarrow d)$, $j_d$ sends a map $(f, \phi) \colon (d_1, x_1) \ra (d_2, x_2)$ with $f_2f = f_1$ to a map $H(f_1)x_1 \ra H(f_2)x_2$ in $H(d)$ defined by $H(f_2)\phi$
\item[(c)]
The adjunction counit is $id \colon j_di_d \Ra 1_{p^{-1}d}$
\item[(d)]
The adjunction unit $1_{(p \downarrow d)} \Ra i_dj_d$ maps 
\begin{center}
$((d_1, x_1), f_1 \colon d_1 \ra d) \Ra ((d, H(f_1)x_1), 1_d \colon d \ra d)$
\end{center}
via the map $(d_1, x_1) \ra (d, H(f_1)x_1)$ in $\int_{\cat D} H$ defined by $(f_1, 1_{H(f_1)x_1})$.
\end{enumerate}
The proof of (2) is dual.
\end{proof}

As a consequence we can prove a Fubini-type formula for the computation of (co)limits indexed by Grothendieck constructions.

\begin{prop}
\label{prop:grothendieckconstrcolims}
\mbox{}
\begin{enumerate}
\item
Suppose that $H \colon {\cat D} \ra {\cat Cat}$ is a functor from a small category to the category of small categories. Let ${\cat M}$ be a category, and $X$ be a $\int_{\cat D} H$-diagram of ${\cat M}$. Assume that the inner colimit of the right side of the equation below exists for all $d \eps {\cat D}$. If either the left side or the outer right side colimits exist, then both exist and we have a natural isomorphism
\begin{center}
$\colim^{\int_{\cat D} H}X \cong \colim^{d \eps {\cat D}} \colim^{H(d)} X$
\end{center}
\item
Suppose that $H \colon {\cat D}^{op} \ra {\cat Cat}$ is a functor from a small category to the category of small categories. Let ${\cat M}$ be a category, and $X$ be a $\int^{\cat D} H$-diagram of ${\cat M}$. Assume that the inner limit of the right side of the equation below exists for all $d \eps {\cat D}$. If either the left side or the outer right side limits exist, then both exist and we have a natural isomorphism
\begin{center}
$\lim^{\int^{\cat D} H}X \cong \lim^{d \eps {\cat D}} \lim^{H(d)^{op}} X$
\end{center}
\end{enumerate}
\end{prop}

\begin{proof}
We only prove (1) - statement (2) uses a dual proof.

For any object $d$ of ${\cat D}$, the inclusion $p^{-1}d \ra (p \downarrow d)$ has a left adjoint, therefore by \lemmaref{lem::adjfunctcofinalfunct} it is homotopy right cofinal, and in particular right cofinal. We identify $p^{-1}d \cong H(d)$. Since $\colim^{H(d)}X$ exists for all $d$, by \propositionref{prop:cofinallimit} $\colim^{(p \downarrow d)}X \cong \colim^{H(d)}X$ exists for all $d$, and by \lemmaref{lem::computerelcolim} $\colim^p X$ exists. By \lemmaref{lem::colimcomposition}, if either $\colim^{d \eps {\cat D}} \colim^{H(d)} X \cong \colim^{\cat D}\colim^p X$ or $\colim^{\int_{\cat D} H}X$ exist, then they both exist and are canonically isomorphic.
\end{proof}

If $H \colon {\cat D}_1$ $\ra$ ${\cat Cat}$ is a constant functor with value ${\cat D}_2$, then the Grothendieck construction $\int_{\cat D} H$ is {\it isomorphic} to the product of categories ${\cat D}_1 \times {\cat D}_2$. In this case, \propositionref{prop:grothendieckconstrcolims} yields


\begin{cor}
\label{cor:grothendieckconstrcolims}
Suppose that ${\cat D}_1$, ${\cat D}_2$ are two small categories and suppose that $X$ is a ${\cat D}_1 \times {\cat D}_2$ diagram in a category ${\cat M}$.
\begin{enumerate}
\item
Assume that the inner colimit of the right side of the equation below exists for all $d_1 \eps {\cat D}_1$. If either the left side or the outer right side colimits exist, then both exist and we have a natural isomorphism
\begin{center}
$\colim^{{\cat D}_1 \times {\cat D}_2}X \cong \colim^{d_1 \eps {\cat D}_1} \colim^{\{d_1\} \times {\cat D}_2} X$
\end{center}
\item
Assume that the inner limit of the right side of the equation below exists for all $d_1 \eps {\cat D}_1$. If either the left side or the outer right side limits exist, then both exist and we have a natural isomorphism
\begin{center}
$\lim^{{\cat D}_1 \times {\cat D}_2}X \cong \lim^{d_1 \eps {\cat D}_1} \lim^{\{d_1\} \times {\cat D}_2} X$ $\square$
\end{center}
\end{enumerate}
\end{cor}

\begin{rem}
\label{rem:grothendieckconstrcolims}
A typical application of \corollaryref{cor:grothendieckconstrcolims} (1) will be for the case when $\colim^{\{d_1\} \times {\cat D}_2} X$ and $\colim^{{\cat D}_1 \times \{d_2\}} X$ exist for all objects $d_1 \eps {\cat D}_1$ and $d_2 \eps {\cat D}_2$. Then in the equation below if either the left or the right outer side colimits exist, they both exist and we have a natural isomorphism
\begin{center}
$\colim^{d_1 \eps {\cat D}_1} \colim^{\{d_1\} \times {\cat D}_2} X \cong \colim^{d_2 \eps {\cat D}_2} \colim^{{\cat D}_1 \times \{d_2\}} X$ 
\end{center}
\end{rem}

